%% LyX 2.4.5 created this file.  For more info, see https://www.lyx.org/.
%% Do not edit unless you really know what you are doing.
\documentclass[11pt,english]{article}
\usepackage[T1]{fontenc}
\usepackage[utf8]{inputenc}
\usepackage{color}
\usepackage{array}
\usepackage{float}
\usepackage{booktabs}
\usepackage{amsmath}
\usepackage{graphicx}
\usepackage{geometry}
\geometry{verbose,tmargin=1.1in,bmargin=1.1in,lmargin=1.1in,rmargin=1.1in}
\usepackage{setspace}
\onehalfspacing
\raggedbottom

\makeatletter

%%%%%%%%%%%%%%%%%%%%%%%%%%%%%% LyX specific LaTeX commands.
%% Because html converters don't know tabularnewline
\providecommand{\tabularnewline}{\\}
%% A simple dot to overcome graphicx limitations
\newcommand{\lyxdot}{.}


%%%%%%%%%%%%%%%%%%%%%%%%%%%%%% User specified LaTeX commands.
%
\usepackage{amsfonts}\newtheorem{prop}{Proposition}
\newtheorem{cor}{Corollary}
\newtheorem{lem}{Lemma}
\interfootnotelinepenalty=10000
\gdef\@makecol{%
   \setbox\@outputbox \box\@cclv
   \let\@elt\relax
   \xdef\@freelist{\@freelist\@midlist}%
   \global \let \@midlist \@empty
   \@combinefloats
   \ifvbox\@kludgeins
     \@makespecialcolbox
   \else
     \setbox\@outputbox \vbox to\@colht {%
       \@texttop
       \dimen@ \dp\@outputbox
       \unvbox \@outputbox
       \vskip -\dimen@
       \ifvoid\footins
         \@textbottom
       \else
         \vfill
         \vskip \skip\footins
         \color@begingroup
           \normalcolor
           \footnoterule
           \unvbox \footins
         \color@endgroup
       \fi
       }%
   \fi
   \global \maxdepth \@maxdepth
}
\setcounter{MaxMatrixCols}{30}
\providecommand{\U}[1]{\protect\rule{.1in}{.1in}}
\newcommand{\rev}[1]{#1}
\newcommand{\revmath}[1]{#1}
\newcommand{\zac}[1]{#1}
\newenvironment{proof}[1][Proof]{\noindent\textbf{#1.} }{\ \rule{0.5em}{0.5em}}

\newtheorem{theorem}{Theorem}
\newtheorem{acknowledgement}{Acknowledgement}
\newtheorem{algorithm}[theorem]{Algorithm}
\newtheorem{axiom}[theorem]{Axiom}
\newtheorem{assumption}{Assumption}
\newtheorem{case}[theorem]{Case}
\newtheorem{claim}[theorem]{Claim}
\newtheorem{conclusion}[theorem]{Conclusion}
\newtheorem{condition}[theorem]{Condition}
\newtheorem{conjecture}[theorem]{Conjecture}
\newtheorem{corollary}[theorem]{Corollary}
\newtheorem{criterion}[theorem]{Criterion}
\newtheorem{definition}{Definition}
\newtheorem{example}[theorem]{Example}
\newtheorem{exercise}[theorem]{Exercise}
\newtheorem{lemma}{Lemma}
\newtheorem{notation}[theorem]{Notation}
\newtheorem{problem}[theorem]{Problem}
\newtheorem{proposition}{Proposition}
\newtheorem{remark}[theorem]{Remark}
\newtheorem{solution}[theorem]{Solution}
\newtheorem{summary}[theorem]{Summary}
\usepackage{lscape}

\makeatother

\usepackage{babel}
\usepackage[hidelinks]{hyperref}
\newcommand{\paperlink}[2]{\hyperref[ref:#1]{#2}}
\newcommand{\bibtarget}[2]{\phantomsection\label{ref:#1}#2}
\newcommand{\figlink}[1]{\hyperref[#1]{Figure~\ref*{#1}}}
\newcommand{\figlinklower}[1]{\hyperref[#1]{figure~\ref*{#1}}}
\newcommand{\eqnlink}[1]{\hyperref[#1]{equation~\ref*{#1}}}
\newcommand{\reflink}[1]{\hyperref[#1]{\ref*{#1}}}
\begin{document}
\title{NIMBYism and the Housing Supply\thanks{We thank St\'ephane Bonhomme, Fatih G\"uvenen, Serdar
\"Ozkan and Roman \v{S}ustek for their comments and help with data
and calibration issues. We also thank participants in seminars at
Durham University, Monash University, University of Sydney, CEMAP,
WAMS, OzMac, SCE and SEHO. All remaining errors are our own.}}
\author{Isaac Gross\thanks{isaac.gross@monash.edu}\\
{\small Monash University} \and David Chivers\thanks{david.chivers@durham.ac.uk\textbf{ }}
\\
{\small Durham University} \\
}
\maketitle
\begin{abstract}
\ We develop a political economy model of the housing market in which
the supply of new homes is shaped by a democratic vote among households.
Voter support for additional housing is determined by the effect this
vote has on households' expected discounted utility via changes in
house prices. Households that own housing, and therefore typically
experience financial losses from additional supply, generally oppose
new housing: the \zac{``Not In My Backyard''} (NIMBY) effect. Households that
wish to buy a house and would benefit from cheaper housing support
more supply. We find that support for additional supply is highest
among households that are younger, low-income or are currently renting.
We then simulate the impact of a temporary baby boom in the population.
We find that those born during the baby boom attain higher rates of
home ownership and greater net worth due to their greater political
influence relative to other generations. However, this comes at the
expense of higher house prices, lower home-ownership rates and less
wealth for the subsequent generations. Calibrating the model to US
demographic data, we find that this political economy mechanism can account
for a 15 per cent increase in house prices since 1950. Using demographic
projections, we project that population ageing will continue to place
upward pressure on house prices into the 21st century.

\ 

\textbf{JEL Classification}: E24, J10, R31. 

\textbf{Keywords}: Housing supply, inequality, political economy
\end{abstract}
\pagebreak{}

\section{Introduction}

The ``Not In My Backyard'' (NIMBY) effect describes the phenomenon
of local residents opposing new housing developments in their neighbourhood.
As the regulation of housing supply is often \zac{decentralised} to local
governments, residents can exert significant political power over
any new developments within their area. The aggregate effect of NIMBYism
can thus negatively impact the housing supply within an economy. In
the US, for example, land is divided into ``zones'' which determine
the size and use of new developments.\footnote{Zoning laws were originally developed in the 1900s in order to improve
``light and air'' in residential areas - see \paperlink{fischel2004}{Fischel (2004)} for
a historical overview. Land use regulation has also been used by white
residents in the US to maintain and increase racial segregation
(\paperlink{sahn2024}{Sahn 2024}; \paperlink{trounstine2020}{Trounstine 2020}).} Only detached, single-family homes are allowed to be built in ``R1''
residential zones, even if high demand for multi-home housing (such
as apartments) exists. Those who own homes within R1 residential areas
have strong incentives to maintain the status quo at the expense of
non-homeowners, who would benefit from multi-home housing in order
to get on the housing ladder (\paperlink{manville_etal2020}{Manville et al. 2020}). 

We develop a political economy model of the housing market in order
to estimate the impact of NIMBYism and show how changing demographic
trends can account for part of the recent rise in house prices. Specifically,
we allow housing supply in a standard life cycle model to be determined
by a democratic process at the local level: households vote according
to their personal preferences and the final outcome is determined
by a median-voter condition.

We then use this model to show that an ageing
society exacerbates the effects of NIMBYism and implies a 15 per cent
increase in house prices since 1950. Using demographic projections,
we project that population ageing will continue to increase political
pressure on the housing market in the 21st century. Finally we examine the dynamics of a temporary
baby boom and \zac{analyse} how it affects the level and distribution of
wealth in the economy. We find that the generation born during the
baby boom shifts the median voter over its life cycle, first toward
additional construction and later toward restrictions on new supply. This
increases its housing wealth and total net worth, but comes at the cost of
decreased home-ownership rates and net worth for subsequent generations.

Our theoretical model follows the standard life cycle housing model
as described by \paperlink{iacoviello2005}{Iacoviello (2005)}. Over the life-cycle, an individual
household typically rents housing until they have saved enough to
meet their borrowing constraint and be able to put a deposit down
on a house. Households also participate in the housing market by either
upsizing (as their wealth grows or positive income shocks occur) or
selling (if they receive a negative income shock or die). Instead
of fixing the housing supply (as in Iacoviello 2005) or having it
produced inelastically by a construction sector (as in Kaplan et al.
2020), we allow the supply of housing to be determined by a democratic
vote of local households. 

Household preferences over the supply of housing are formed according
to the impact that a small change in price will have on their expected
discounted utility. Households who gain from marginally higher prices
will vote against new developments, whereas those who marginally lose
out from higher house prices will vote in favour of new developments.\footnote{Households assume that increasing the housing supply decreases house
prices within their local area. Although increasing housing supply
may increase house prices via agglomeration effects, these are dispersed
across a larger area (\paperlink{duranton_kerr2015}{Duranton and Kerr, 2015}) and so will be ignored
by households voting within decentralised local government authorities
that make up a greater metropolitan area. This dynamic underpins the
NIMBY position: in principle, they may support increased housing supply
- just not in their own backyard. } NIMBYist views in our environment are thus largely determined by households'
ownership of housing and their expected probability of being a net-buyer
in future periods. NIMBYist households are disproportionately older,
wealthier, have higher incomes and are more likely to be homeowners,
which is in line with the empirical literature surrounding NIMBYism
(see \paperlink{whittemore_bendor2018}{Whittemore \& BenDor (2018)}, \paperlink{scally_tighe2015}{Scally and Tighe (2015)}, and \paperlink{hankinson2018}{Hankinson (2018)} for a range of surveys on the subject). 

Our model primarily focuses on the financial incentives of NIMBYism,
as households' own homes are the primary asset for most households.\footnote{Excluding the top 1\% of the wealth distribution, home equity makes
up 28.9\% of total US household wealth and is comparable only to retirement
accounts 32.8\%. The next largest asset, stocks and mutual funds,
makes up only 10.2\% of total household wealth (see Eggleston et al.
2020).} However there are several non-financial motivations that may lead
a household to oppose new developments - ranging from infrastructure
capacity concerns to simply disliking the proposed style of architecture
or neighbourhood characteristics (\paperlink{ahlfeldt2011}{Ahlfeldt 2011}, \paperlink{bosetti_sims2016}{Bosetti and Sims 2016})\textit{.}
We allow for these non-financial preferences to affect households'
votes but assume these preferences are stable over the life cycle\textit{.}
The financial incentive for NIMBYism, however, does vary over the
life cycle and is analogous to an insider-outsider problem: outsiders
(renters) would prefer more houses to be built until they become an
insider (homeowner), and then seek to restrict supply. 

As home-ownership rates and thus support for new developments change
over the life cycle, the demographic make-up of an economy will play
an important role in the aggregate effect of NIMBYism. During the
mid 1940s, a baby boom occurred in a number of OECD countries including
the US, UK, New Zealand, Australia, France and Canada (\paperlink{greenwood_etal2005}{Greenwood et al. 2005}). This baby boom lasted until the subsequent baby bust of
the mid 1960s when birth rates returned to trend. The effect of the
spike in the birth rate led to a generation known as the Baby Boomers
which has taken up a significant share of the population over time.
As those who are close in age will have similar preferences over housing
on average (as a result of being at a similar stage of the life-cycle),
those born during the baby boom will have outsized political power
and thus have a greater influence over the aggregate housing stock.
This is perhaps one reason why rents are increasing relative to median
income over the long run as shown in \figlink{fig:1}. It is also why the
US home-ownership rate has been falling for younger people over time as shown
in \figlink{fig:2}.
\begin{figure}[H]
\centering
\begin{minipage}[t]{0.48\textwidth}
\centering
\includegraphics[width=\linewidth]{Figures/figure1_rent_to_median_income_vector.pdf}
\caption{Rent to Median Income}
\label{fig:1}
\vspace{0.25em}
{\scriptsize Data source: CPS (2022).}
\end{minipage}\hfill
\begin{minipage}[t]{0.48\textwidth}
\centering
\includegraphics[width=\linewidth]{Figures/figure2_home_ownership_rate_vector.pdf}
\caption{Home Ownership by Age Cohort}
\label{fig:2}
\vspace{0.25em}
{\scriptsize Data source: CPS (2022).}
\end{minipage}
\end{figure}

In this paper, we estimate how the changing demographic shares across
ageing economies have changed the political-economic equilibrium in
the housing market. As economies age, the balance of political power
shifts in favour of NIMBYist households resulting in lower housing
supply and higher house prices. We find that house prices increased
by around 15 per cent due to these political-economic effects since
the 1950s. Using demographic projections, as shown in \figlinklower{fig:3},
we find that the ageing economy is expected to keep the house price
at or above its current level until 2100.

Finally, we use our political-economic model of the housing market
to examine the distributional consequences of recent demographic shocks.
We find that a temporary boom in the birth rate will initially decrease
inter- and intra-generational wealth inequality as the boom generation
initially supports lower house prices when they are disproportionately
young renters. As the cohort ages, this effect reverses and the cohort
exerts greater political pressure against the supply of new housing. This
decrease in support for housing supply, and the resulting higher house
price, leads to a persistent scarring effect on future generations
who are unable to achieve home ownership at the same rate as previous
generations, even after the initial baby boom generation has completely
died off. 
\begin{figure}[H]
\begin{centering}
\includegraphics[width=0.68\textwidth]{Figures/figure3_age_share_projection_vector.pdf}
\par\end{centering}
\noindent\raggedright{}\caption{Age Share of Population Over Time}
\label{fig:3}
\vspace{0.35em}
\begin{centering}
\begin{minipage}{0.70\textwidth}
{\scriptsize Data source: UN (2019). Note: dashed line indicates population
forecasts from 2020 onward assuming medium immigration.}{\scriptsize\par}
\end{minipage}
\par\end{centering}
\end{figure}
 
The remainder of the paper proceeds as follows. Section 2 discusses
the related literature. Section 3 presents the model and voting mechanism.
Section 4 describes the calibration. Section 5 reports the stationary
and transition exercises. Section 6 presents robustness exercises, and
Section 7 concludes.


\section{Related Literature }

This paper is related to the following strands of the literature.
Firstly, there is growing literature surrounding the impact of regulation
on the housing supply. The concept of NIMBYism is often attributed
to \paperlink{frieden1979}{Frieden (1979)}, who noted a rise in environmental regulations to
impede new housing developments after the 1970s. \paperlink{gyourko_mayer_sinai2013}{Gyourko et al. (2013)}
support this view and argue that before the 1970s, housing construction
tended to follow areas of high price growth. However, establishing
a direct link between regulation and the housing supply is difficult
due to the idiosyncratic forms that regulations can take in different
localities (\paperlink{gyourko_molloy2014}{Gyourko and Molloy 2014}). To get round this, Glaeser et
al. (2006) offer indirect evidence of excess regulation by noting
the existence of a relatively large wedge between prices and marginal
construction costs. The difficulty in providing direct evidence stems
from the endogeneity between land use regulation and the housing market.
Recent papers have used boundary discontinuity design to address this
issue (\paperlink{kahn_vaughn_zasloff2010}{Kahn et al. 2010}, \paperlink{turner_etal2014}{Turner et al. 2014}, \paperlink{kulka_etal2022}{Kulka et al. 2022}). For example,
\paperlink{anagol_etal2021}{Anagol et al. (2021)} exploit a zoning reform in Sao Paulo and show
that housing developers apply for more permits in blocks with higher
maximum densities. In their equilibrium model of housing supply this
leads to a 2.2\% increase in the total housing stock and 0.5\% average
price reduction. Similarly, \paperlink{song2021}{Song (2021)} constructs a nationwide data
set of minimum lot area estimates and finds that doubling the US
minimum lot area size would increase house prices by 14\% and rents
by 6\%.

In theory, an increase in the supply of housing units may increase
house prices via improved amenities and neighbourhood quality, or
via agglomeration effects which increase economic activity (\paperlink{herkenhoff_ohanian_prescott2018}{Herkenhoff et al. 2018}, \paperlink{chatterjee_eyigungor2017}{Chatterjee and Eyigungor 2017}, \paperlink{titman_zhu2024}{Titman and Zhu 2024}).
In the immediate local vicinity, however, these demand effects do
not appear to dominate the supply channel. For example, \paperlink{li2022}{Li (2022)}
finds that within 500 ft, for every 10\% increase in the housing stock,
rents decrease by 1\% and for every 10\% increase in the condo stock,
condo sales prices decrease by 0.9\%. Similarly, \paperlink{pennington2021}{Pennington (2021)}
finds that rents fall by about 2\% for properties within 100 metres
of a new housing project, and the effects decrease to zero when 1.5
km away. This illustrates the coordination problem in NIMBYism: the impact of extra housing within
a homeowner's immediate vicinity will negatively affect prices, whereas
the benefits of agglomeration within the wider area can only occur if
local residents accept extra housing in their back yard. As a result
an emerging political economy literature has started to address this
issue (see for example, \paperlink{parkhomenko2023}{Parkhomenko 2023}, \paperlink{bunten2017}{Bunten 2017}, \paperlink{duranton_puga2023}{Duranton and Puga 2023} and \paperlink{albouy_etal2019}{Albouy et al., 2019}). Using a difference-in-differences
approach, \paperlink{mast2024}{Mast (2024)} shows that when elections of town councils were
localised from \textquotedblleft at-large\textquotedblright{}
to \textquotedblleft ward\textquotedblright{} areas, housing permits
decreased by 20\%.

This paper is also related to an older strand of literature on political
redistribution, such as \paperlink{alesina_rodrik1994}{Alesina and Rodrik (1994)} and \paperlink{persson_tabellini1994}{Persson and Tabellini (1994)}. The main idea of these papers is that when income
and resources are distributed heterogeneously, the majority may vote
for redistributive policies. These papers rely on a voting mechanism,
the median-voter theorem, to determine what policy will be pursued.
The theorem shows that when voting takes place over a single issue,
and preferences are single-peaked, then the policy chosen under a
democracy will be the one chosen by the median-voter. 

Housing politics is a good candidate for the median-voter theorem
as voting takes place at the local level and usually on specific propositions
related to relaxing zoning regulations. The idea of homeowners as
a specific voting class in local government decisions as a result of
financial interests is related to the ``homevoter hypothesis'' (\paperlink{fischel1999}{Fischel 1999}; \paperlink{fischel2001}{Fischel 2001}). \paperlink{einstein_etal2019}{Einstein et al. (2019)} examine planning meetings in Massachusetts
and find that homeowners and older residents are overly representative
in planning meetings. This is in line with voter turnout in local
elections which in general tend to be from these two groups (\paperlink{kogan_etal2018}{Kogan
et al. 2018}, \paperlink{oliver_ha2007}{Oliver and Ha 2007}). Similarly, \paperlink{jiang2018}{Jiang (2018)} finds that
voter turnout increases for mayoral elections, relative to presidential
elections, when home-ownership rates increase. In addition, \paperlink{cheung_meltzer2013}{Cheung and Meltzer (2013)} find that the prevalence of Homeowners Associations
is positively associated with an increase in land use regulations.
Similarly, \paperlink{hilber_robert_nicoud2013}{Hilber and Robert-Nicoud (2013)} find a positive correlation
between the degree of physical development in an area and land use
regulation. Both \paperlink{ahlfeldt_maennig2015}{Ahlfeldt and Maennig (2015)} as well as \paperlink{dehring_depken_ward2008}{Dehring and Ward (2008)} use public referendums to directly test the homevoter
hypothesis. \paperlink{ahlfeldt_maennig2015}{Ahlfeldt and Maennig (2015)} use a referendum on the expansion
of an airport in Berlin as a natural experiment and find greater support
among homeowners as opposed to leaseholders. The support for the proposal
comes from two effects: a consumption effect (through use of the amenity)
as well as a wealth effect (via house price increases). \paperlink{dehring_depken_ward2008}{Dehring and Ward (2008)} use pre-vote changes in residential property values to
test the probability of support for the new stadium for the Dallas
Cowboys and find between a 0.9\% and 1.2\% increase in support for
every \$1000 dollar increase in house prices.

Numerous surveys of voters have attempted to describe the demographics
of voters who oppose new housing supply. \paperlink{whittemore_bendor2018}{Whittemore and BenDor (2018)}
find that older, less educated were more likely to oppose new developments,
though they found no significant relationship between home ownership
and NIMBYist views. By contrast \paperlink{hankinson2018}{Hankinson (2018)} and \paperlink{scally_tighe2015}{Scally and Tighe (2015)} found that home owners were significantly more likely to oppose
new housing supply - particularly in their locality.

Finally, this paper is related to the literature surrounding life-cycle
models with housing markets. Generally, life-cycle models assume exogenous
house prices due to computational complexity (\paperlink{fernandez_krueger2011}{Fernandez-Villaverde and Krueger (2011)}, \paperlink{iacoviello_pavan2013}{Iacoviello and Pavan (2013)}, \paperlink{yang2009}{Yang (2009)}). There
are, however, relatively few papers where aggregate shocks affect
house prices (\paperlink{justiniano_etal2015}{Justiniano et al. 2015}). Our paper is most closely related
to \paperlink{ortalo_magne_prat2014}{Ortalo-Magn\'e and Prat (2014)}, who use a median-voter theorem in
an overlapping\textendash generations economy to show that increasing
home-ownership rates decrease housing affordability. Similarly, \paperlink{mankiw_weil1989}{Mankiw and Weil (1989)} predict the growth rate of housing demand via changes
in the age distribution alone and more recently \paperlink{mumtaz_sustek2023}{Mumtaz and \v{S}ustek (2023)} find similar results except housing demand peaks in the 2000s.

\section{The Model}

We outline a life-cycle model with housing and debt
choices. The model features heterogeneous, life-cycle households that
\zac{optimise} non-durable consumption, housing and
saving in liquid assets. A number of factors influence households'
political support for additional housing supply. These include households'
income and wealth, portfolio allocation, housing tenure and age over the
life cycle. As demographics change, the composition of the voting population
changes, and this alters political support for new construction. We use the
model to study how an exogenous demographic shock changes voting over housing
supply and, through this political channel, equilibrium house prices.

\rev{Households make housing and saving decisions in an environment in which
future housing-supply restrictions are chosen politically. A demographic shock
changes the future age and tenure composition of voters. When a larger cohort
is young, it raises the relative weight of renters and prospective buyers, who
benefit from additional construction. As the same cohort ages into
home ownership, political support shifts toward tighter restrictions on new
supply. \zac{These changes in voter composition affect future supply restrictions and the market-clearing house prices associated with permitted supply.}}

\rev{Households are forward-looking, so the dynamic problem requires
expectations about future housing-market conditions. These expectations enter
through continuation values, the value of reaching future periods with a given
state. We specify expectations using forecast rules that differ in the
information households use about how future demographic change affects voting
behaviour and, through the political process, future restrictions on housing
supply. One forecast rule
uses the current housing-market environment to \zac{summarise} future conditions:
households do not explicitly forecast the future demographic path or the
political outcomes it will generate. The alternative forecast rules use the
demographic path and the political housing-supply mechanism to forecast future
voter composition, supply restrictions and prices. Comparing these forecast
rules isolates how expectations about future demographics and politics shape
\zac{capitalisation}, housing choices and voting incentives.}

A consumer of age $j$ enters the economy at age 25 and lives for
\rev{$J$ \zac{periods}, indexed by $j=0,\ldots,J-1$, and \zac{maximises}}
lifetime expected utility which
consists of consumption $c_{t}$, housing services, $h_{t}$ and a
\rev{terminal} bequest \rev{$v\left(w_{J}\right)$}. Housing services depend on whether
they own or rent the property. Individuals also have access to a liquid
asset that pays an interest $r_{t+1}$. The household's lifetime utility
function is given by

\begin{equation}
\revmath{\sum^{J-1}_{j=0}\beta^{j}u(c_{j},h_{j})+\beta^{J}v\left(w_{J}\right)}
\end{equation}
where households \zac{optimise} over consumption, $c_{t}$, liquid assets,
$a_{t}$, and housing, $h_{t}$ at each age $j$. The flow utility
is increasing and strictly concave in $c$ and $h$. Households receive
utility from wealth they bequeath \rev{$w_{J}$}. We assume the functional
form 

\begin{equation}
u(c_{t},h_{t})=\frac{\left(c^{1-\zeta}_{t}h^{\zeta}_{t}\right)^{1-\gamma}}{1-\gamma}
\end{equation}
where the consumption of housing stock consists of either the amount
of housing owned or rented

\begin{equation}
h_{t}=\begin{cases}
\theta^{rent}h^{rent}_{t} & \mbox{if }h^{own}_{t}=0\\
h^{own}_{t} & \mbox{if }h^{own}_{t}>0.
\end{cases}
\end{equation}
The curvature parameter $\gamma$ determines households' risk aversion
and inter-temporal elasticity of substitution and $\zeta$ denotes
housing's share of total consumption. Rental housing provides a stream
of utility that is discounted by $\theta^{rent}$ compared to self-owned
housing. Finally there is a warm glow bequest function when households
die that provides utility
\begin{equation}
v(w)=\psi\frac{\left(w+\overline{w}\right)^{1-\gamma}}{1-\gamma}
\end{equation}
which generates a bequest function with decreasing marginal utility
with respect to wealth, $w$, governed by the parameters $\overline{w}$
and $\psi$. Households supply one unit of labour inelastically and
are paid a wage $y_{t,j}$. The wage they receive varies according
to their productivity which is driven by the exogenous process

\begin{eqnarray}
\revmath{\log y_{i,j,t}} & \revmath{=} & \revmath{g_{j}+z_{i,t}+\varepsilon_{i,t}\text{ for }j<J^{retired},}\nonumber \\
\revmath{z_{i,t+1}} & \revmath{=} & \revmath{\rho z_{i,t}+\eta_{i,t+1}.}
\end{eqnarray}
\rev{The deterministic component $g_{j}$ is the age profile. The persistent
component $z_{i,t}$ and transitory component $\varepsilon_{i,t}$ are
approximated with a finite-state Markov process calibrated to yield realistic
income inequality, following \paperlink{kindermann_krueger2022}{Kindermann and Krueger (2022)}. The numerical
process also includes two top-income states that capture ``superstar''
income realizations. Once workers reach retirement age $J^{retired}$,
their retirement income is fixed as a function of income at retirement and is
then carried as a retirement-income state.} The income process is \zac{normalised} to have a mean of
1, such that the price of housing $p^{h}_{t}$ reflects the house
price to income ratio.

Households may save or borrow liquid assets subject to a borrowing
constraint where each household can borrow against a proportion, $m$,
of the total value of the stock of housing they own

\begin{equation}
a_{t}\geq-mh^{own}_{t}p^{h}_{t}.
\end{equation}
\rev{Net asset positions earn the risk-free rate when \zac{non-negative} and are
financed at the mortgage rate when negative. Let}
\begin{equation}
\revmath{r_{t}\left(a_{t}\right)=
\begin{cases}
r^{a}_{t} & \mbox{if }a_{t}\geq0,\\
r^{m}_{t} & \mbox{if }a_{t}<0,
\end{cases}}
\end{equation}
\rev{where $r^{m}_{t}=r^{a}_{t}+\chi$ and $\chi$ is the mortgage spread.}
Households may buy or sell housing at price, $p^{h}_{t},$ subject
to a minimum block size $\bar{h}$. \rev{Thus owned housing satisfies
$h_{t}^{own}\in\{0\}\cup[\bar h,\infty)$.} All property sales, including
those that are sold as part of a bequest are subject to a transaction
cost $F_{sell}$ proportional to the total value of property sold,
which is reflected by the indicator function $I_{sell}$. If they
do not own any housing they must rent it from real estate firms at
price $p^{rent}_{t}$. The household's utility function is thus \zac{maximised}
subject to the general budget constraint\footnote{Noting that households only interact with the rental market if they
own no housing stock of their own.} 

\begin{equation}
a_{t+1}=\left(1+r_{t}\left(a_{t}\right)\right)a_{t}+y_{t}-c_{t}-p^{rent}_{t}h^{rent}_{t}-\left(1-I^{sell}F^{sell}\right)p^{h}_{t}\varDelta h^{own}_{t}
\end{equation}
\rev{where $I^{sell}_{t}=\mathbf{1}\{\varDelta h^{own}_{t}<0\}$ is an indicator
for households that choose to sell part or all of their current stock of
housing.}

\rev{New households enter with initial liquid wealth drawn from the calibrated
entrant-wealth distribution. The mean of this distribution, $a_{25}$, is tied
to average terminal bequests in the stationary economy.}

\subsection{Household's Problem}

Each household's state vector is $\mathbf{s}_{t}=\left\{ a_{t},h^{own}_{t-1},z_{t}\right\} $,
where $a_{t}$ is the stock of liquid assets, $h^{own}_{t-1}$ is
the amount of housing they owned at the end of period $t-1$ and $z_{t}$ 
is the persistent component of their income. \rev{After retirement, the fixed
retirement-income state described above is also carried forward, but we
suppress it in the notation.} The price of housing,
$p_{t}\equiv p_{t}^{h}$, is an aggregate state variable. \rev{Expectations enter the household
problem through the future price path used in continuation values. Forecast
rules are indexed by $\lambda$. At date $t$, households evaluate housing and
saving decisions using the perceived continuation price path}
\begin{equation}
\revmath{\widehat{\mathbf{p}}_{t}\left(\lambda\right)
=\left(p_{t},\widehat{p}_{t+1|t}\left(\lambda\right),\ldots,\widehat{p}_{T|t}\left(\lambda\right)\right),}
\label{eq:forecastpath}
\end{equation}
\rev{where $\widehat{p}_{t+s|t}\left(\lambda\right)$ is the house price that
households use for date $t+s$. In the numerical \zac{transition path}, the full perceived
object is the collection of these date-specific continuation paths. The
transition value function is therefore $V^{\lambda}_{j,t}\left(\mathbf{s}_{t};
\widehat{\mathbf{p}}_{t}\left(\lambda\right)\right)$. In the Bellman equations
below, $\widehat{p}_{t+1|t}\left(\lambda\right)$ denotes the relevant remaining
continuation path from date $t+1$ onward, written compactly to keep the notation
readable. When $\lambda=0$, the forecast path collapses to
the current-price rule:}
\begin{equation}
\revmath{\widehat{p}_{t+s|t}\left(0\right)=p_{t}\quad\text{for all }s\geq1.}
\label{eq:current_price_forecast}
\end{equation}
\rev{In this case, households do not separately forecast future demographic
changes or the political housing-supply restrictions those changes induce.
Future house prices are therefore assumed to equal the current house price at
date $t$. For larger values of $\lambda$, household forecasts place more
weight on the demographic-political path implied by the model.}

The value function for the household is given by the maximum of the three possible housing tenure
options: rent (either because they owned no housing or sell it all
that period), or are owners who choose to adjust their stock of housing
or maintain their current level of housing. Households who choose
to rent (and thus sell whatever stock of housing they owned at the
beginning of the period) are able to freely adjust the amount of housing
they rent each period. The value function for a household that choose
to rent housing (and thus sell any previously owned housing stock)
at age $j$ is 

\begin{eqnarray}
V^{rent,\lambda}_{j}\left(\mathbf{s}_{t};p_{t}\right) & = & \underset{c_{t},a_{t+1},h^{rent}_{t}}{\mbox{max}}\revmath{u\left(c_{t},\theta^{rent}h^{rent}_{t}\right)}+\beta E_{t}V^{\lambda}_{j+1}\left(\mathbf{s}_{t+1};\widehat{p}_{t+1|t}\left(\lambda\right)\right)\\
s.t\nonumber \\
a_{t+1} & = & \left(1+r_{t}\left(a_{t}\right)\right)a_{t}+y_{t}-c_{t}-p^{rent}_{t}h^{rent}_{t}+\left(1-F^{sell}\right)p^{h}_{t}h^{own}_{t-1}\nonumber \\
a_{t+1} & \geq & 0\nonumber \\
h^{own}_{t} & = & 0\nonumber 
\end{eqnarray}
The value function for home owners who choose not to adjust their
(non-zero) stock of housing in period $t$ is

\begin{eqnarray}
V^{non-adjust,\lambda}_{j}\left(\mathbf{s}_{t};p_{t}\right) & = & \underset{c_{t},a_{t+1}}{\mbox{max}}u\left(c_{t},h^{own}_{t}\right)+\beta E_{t}V^{\lambda}_{j+1}\left(\mathbf{s}_{t+1};\widehat{p}_{t+1|t}\left(\lambda\right)\right)\\
s.t\nonumber \\
a_{t+1} & = & \left(1+r_{t}\left(a_{t}\right)\right)a_{t}+y_{t}-c_{t}\nonumber \\
a_{t+1} & \geq & -mp^{h}_{t}h^{own}_{t}\nonumber \\
h^{own}_{t} & = & h^{own}_{t-1}\nonumber 
\end{eqnarray}
finally, the value function for households who choose to buy or sell
housing in period $t$ is

\begin{eqnarray}
V^{adjust,\lambda}_{j}\left(\mathbf{s}_{t};p_{t}\right) & = & \underset{c_{t},a_{t+1},h^{own}_{t}}{\mbox{max}}u\left(c_{t},h^{own}_{t}\right)+\beta E_{t}V^{\lambda}_{j+1}\left(\mathbf{s}_{t+1};\widehat{p}_{t+1|t}\left(\lambda\right)\right)\\
s.t\nonumber \\
a_{t+1} & = & \left(1+r_{t}\left(a_{t}\right)\right)a_{t}+y_{t}-c_{t}-\left(1-I^{sell}F^{sell}\right)p^{h}_{t}\varDelta h^{own}_{t}\nonumber \\
a_{t+1} & \geq & -mp^{h}_{t}h^{own}_{t}\nonumber \\
h^{own}_{t} & = & h^{own}_{t-1}+\varDelta h^{own}_{t}\nonumber 
\end{eqnarray}
The final value function is the maximum of each of the three choices
over housing tenure

\begin{equation}
V^{\lambda}_{j}\left(\mathbf{s}_{t};p_{t}\right)=\underset{}{\mbox{max}\left\{ V^{rent,\lambda}_{j}\text{\ensuremath{\left(\mathbf{s}_{t};p_{t}\right)}},V^{adjust,\lambda}_{j}\left(\mathbf{s}_{t};p_{t}\right),V^{non-adjust,\lambda}_{j}\left(\mathbf{s}_{t};p_{t}\right)\right\} }
\end{equation}

The expected value function for the final year of life is the same
for all types of housing tenure and is given by the warm glow bequest
function above, where total wealth is the sum of liquid
and illiquid wealth given the current housing price.

\begin{equation}
w=a+\left(1-F^{sell}\right)p^{h}_{t}h^{own}
\end{equation}


\subsection{Real Estate Firms}

Real estate firms borrow from households by issuing debt in the form
of the liquid asset to fund the purchase of housing, which is rented
out to households that do not own their own home. Real estate firms
operate in a perfectly competitive market, and thus the rental price
is pinned down by the cost of borrowing, a fixed cost of providing
rental services, the price of purchasing house stock and any expected
capital appreciation driven by the drift in the price of housing.
We assume that real estate firms are able to borrow against the full
value of their housing stock. The zero profit condition pins down
the price of rental stock as the function of the cost of housing,
capital gains from housing and the interest rate

\begin{equation}
\revmath{p^{rent}_{t}=\phi^{rent}+r^{a}_{t}p^{h}_{t}-\left(\widehat{p}_{t+1|t}\left(\lambda\right)-p^{h}_{t}\right),}
\end{equation}
where $\phi^{rent}$ is the per unit cost of providing rental services.
\rev{The final term is the expected capital gain from holding housing, evaluated
using the same one-period price forecast that enters household continuation
values. Rental firms break even under this perceived price path; ex post
profits can differ from zero when perceived and \zac{realised} price changes differ.
In the numerical exercises we restrict attention to paths for which the implied
rental price remains positive.}

\subsection{Housing Supply}

The aggregate economy consists of a unit measure of neighbourhoods across
which all types of households are equally distributed. \rev{Aggregate demographic
changes are therefore interpreted as common changes in the demographic
composition of representative local housing markets. This abstracts from
spatial sorting and from differences in the local concentration of particular
demographic groups.} The neighbourhood's
housing supply in period $t+1$ is determined in each neighbourhood
by a simple majoritarian voting mechanism in which households vote
over the amount of permitted housing in an area. Only households who
currently reside in that neighbourhood can vote to determine the next
period's level of new housing. Once the quantity of housing permits
has been determined by the households' collective vote within the
neighbourhood, the supply of housing is costlessly produced by a perfectly
competitive construction sector in $t+1$. This is similar to Kaplan
et al. (2020), who model a perfectly competitive construction sector
in which all rents from land permits accrue to the government. \rev{New
construction is physically costless, but permits are scarce. House prices
therefore reflect the scarcity value of permitted housing. Permit rents accrue
to the government or absentee landowners and are not rebated to households.}

An increase in the housing supply will reduce the price of housing
in the locality and households vote understanding that an increase
in housing permits will lead to a small persistent decrease in the
cost of housing. An increase in the housing supply will also affect
the price of rental properties by lowering the input price for the
real estate firms. As a result, each household in a neighbourhood will
vote to increase (or decrease) the availability of housing permits
in their locality based on how the induced change in housing costs affects
their expected welfare. \rev{\zac{Given the quantity of permitted housing, the house price adjusts to clear the local housing market.} We therefore represent the welfare effect of a
marginal supply restriction by the fixed proportional price perturbation
$\epsilon^{p}$ that such a restriction implies. In the transition calculations,
the perturbation is applied to the perceived price path that represents this
persistent price effect. Let
$\widehat{\mathbf{p}}_{t}^{+\epsilon}\left(\lambda\right)$ denote the path that
raises the current and future prices relevant for this marginal restriction by
the factor $1+\epsilon^{p}$, and define}

\begin{equation}
\revmath{\Delta V^{\lambda}_{j,t}\left(\mathbf{s}_{t}\right)
=V^{\lambda}_{j,t}\left(\mathbf{s}_{t};\widehat{\mathbf{p}}_{t}^{+\epsilon}\left(\lambda\right)\right)
-V^{\lambda}_{j,t}\left(\mathbf{s}_{t};\widehat{\mathbf{p}}_{t}\left(\lambda\right)\right).}
\end{equation}
\rev{In the stationary problem this collapses to the corresponding fixed
perturbation of the stationary house price. Households vote to restrict housing
supply when}
\begin{equation}
\revmath{\Delta V^{\lambda}_{j,t}\left(\mathbf{s}_{t}\right)+\sigma\geq0.}
\end{equation}
\rev{The perturbation size $\epsilon^{p}$ is fixed throughout the computations,
so $\sigma$ is calibrated in the same utility units as this finite-difference
welfare comparison.} The parameter $\sigma$ represents a non-financial preference over the supply
of housing. A positive value of $\sigma$ shifts households toward
restricting new supply, while a negative value shifts them toward allowing
additional construction. For example, households may view new housing
negatively as a result of congestion effects on infrastructure, or
positively because of the new amenities that greater density could bring.
Households may also hold ideological views about housing supply, which this
parameter would capture as well.

In principle these preferences could vary by household, either due
to changes in age, productivity, or as a result of an idiosyncratic
shock. However, allowing for such heterogeneity would require introducing
a larger set of additional parameters that would then need to be calibrated.
To limit our degrees of freedom, we therefore calibrate these preferences
as a simple, fixed constant across all households regardless of age,
income or wealth. As a result, this parameter is effectively the average
non-financial preference for all households. We test this assumption
in Section 6 and find that the main results are little changed under
alternative assumptions.

The impact of a small perturbation to the price of housing will vary
depending on household's ownership of liquid assets, housing and their
current level of productivity. If a small increase in the price of
housing increases their expected lifetime utility, then this expression
will be positive and they will vote to restrict housing permits.

\subsection{The Median Voter Theorem}

\zac{Applying the median-voter logic, equilibrium in the housing market occurs
when the median voter is indifferent to a marginal change in housing permits,
evaluated through the price change that permits imply.}\footnote{\zac{Strictly
speaking, the standard median-voter theorem assumes single-peaked preferences
over a one-dimensional policy space. In the quantitative exercise, the relevant
object is the median-voter indifference condition for a marginal change in
permits: at the equilibrium, one half of households support a marginal
restriction on new permits and one half do not.}} \zac{Let $\mu_{j,t}\left(d\mathbf{s}\right)$ denote the distribution
of individual states conditional on age $j$, and let $\omega_{j,t}$ denote the
population weight on that age. For a perceived price path
$\widehat{\mathbf{p}}_{t}\left(\lambda\right)$, the age-weighted vote share for
restricting housing supply is}
\begin{equation}
\revmath{\zac{S_t\left(\widehat{\mathbf{p}}_{t}\left(\lambda\right)\right)
=\sum_{j=0}^{J-1}\omega_{j,t}
\int \Theta_{j,t}\left(\mathbf{s};\widehat{\mathbf{p}}_{t}\left(\lambda\right)\right)
\mu_{j,t}\left(d\mathbf{s}\right),}}
\end{equation}
\zac{where}
\begin{equation}
\revmath{\zac{\Theta_{j,t}\left(\mathbf{s};\widehat{\mathbf{p}}_{t}\left(\lambda\right)\right)
=\mathbf{1}_{\left\{\Delta V^{\lambda}_{j,t}\left(\mathbf{s}\right)+\sigma\geq0\right\}}.}}
\end{equation}
\zac{In the stationary accounting exercise, the perceived path is constant at
$p^{h}$, and the local political equilibrium chooses the house price such that}
\begin{equation}
\revmath{\zac{S_t\left(p^{h}\right)=\frac{1}{2}.}}
\label{eq:median_voter_condition}
\end{equation}
\zac{At this price, one half of the electorate supports a marginal restriction
on new permits and one half does not. We interpret the resulting price as the
market-clearing house price conditional on the permit scarcity implied by the
political outcome. Along a transition, the same vote-share object is applied
period by period to generate the price path.}

The welfare effect of a marginal supply restriction is determined by the
household's current portfolio, tenure status and expected future housing
demand. \figlink{fig:4} shows how these household characteristics
correlate with the decision to restrict supply. Renters and households that
expect to buy additional housing tend to support additional permits because
new supply lowers purchase and rental costs. These households tend to be
younger, have lower incomes and hold less wealth. By contrast, homeowners
who are unlikely to buy additional housing tend to support restrictions
because lower supply preserves the resale value of their housing wealth and
the size of their bequest.

The induced equilibrium price of housing therefore varies with the relative
size of these groups. An exogenous increase in young households, who
disproportionately rent and expect to buy housing in the future, increases
political support for additional supply. The house price then falls until
the median household is indifferent over a marginal change in housing
permits.

In order to abstract from strategic voting across dynasties, we assume
different children and parents live in different neighbourhoods and
thus do not affect the other's housing-market dynamics.\footnote{According to \paperlink{fu2019}{Fu (2019)} 61.5\% adults live more than 10 miles away
from their mother.} Consequently even if an individual may benefit from a larger inheritance
if the price of their parents home increases, it will not affect their
vote on the optimal level of housing supply in their local neighbourhood.\footnote{This assumption also makes our model robust to other forms of altruistic
giving. For example, in addition to receiving a warm glow from giving
to their descendants at death, households may care about their descendant's
welfare directly and thus may wish to sell housing in order to give
their children funds to purchase housing earlier. As families live
in different neighbourhoods, however, this would still not change
their voting behaviour. This is because parents with housing would
still wish to restrict housing supply in their neighbourhood in order
to receive the highest return on their house so they can maximise
their child's welfare.}

We also assume that households are unable to form coalitions when voting,
nor bargain political support across time periods and vote solely
based on their individual interest at time $t$.\footnote{This assumption allows us to abstract from any potential employment
benefits derived from additional demand for the construction sector. }

\begin{figure}[H]
\begin{centering}
\includegraphics[width=0.82\textwidth]{Figures/figure4_unified_model_mechanism.pdf}
\par\end{centering}
\noindent\raggedright{}\caption{Household Support for Restricting Housing Supply}
\emph{\label{fig:4}}
\end{figure}


\subsection{Stationary Equilibrium}

For ease of notation, define the age distribution as $\Omega=\{\omega_{25},...\omega_{80}\}$,
where $\omega_{i}$ is the share of the population at age $i$.\footnote{We drop the $t$ subscript for simplicity.}
\rev{For each exogenous age distribution $\Omega$, we solve a stationary
political-economic exercise: we hold age weights fixed and ask what house price
satisfies the median-voter condition. The age distribution is therefore taken
from the data or projection, rather than generated inside the model.} The stationary accounting equilibrium
consists of value functions, decision rules over consumption,
$c\left(\mathbf{s};p^{h}\left(\Omega\right)\right)$,
housing demand, $\left\{ h^{rent}\left(\mathbf{s};p^{h}\left(\Omega\right)\right),h\left(\mathbf{s};p^{h}\left(\Omega\right)\right)\right\} $,
liquid savings $a\left(\mathbf{s};p^{h}\left(\Omega\right)\right)$,
and votes over housing supply, \rev{$\left\{\Theta_{j}\left(\mathbf{s};p^{h}\left(\Omega\right)\right)\right\}_{j=0}^{J-1}$},
which define the households' optimal allocations as a function of
their individual state variables and the aggregate price of housing,
and the prices of owned and rental housing, $p^{h}\left(\Omega\right)$
and $p^{rent}\left(\Omega\right)$. \rev{Within each accounting exercise,
household formation is chosen only to hold the imposed demographic structure
fixed.} \rev{The transition calculations retain the same household decision
problem, replacing the stationary house price with the forecast price path
defined above.} \rev{The age-specific distributions
$\left\{\mu_{j}\left(d\mathbf{s};p^{h}\left(\Omega\right)\right)\right\}_{j=0}^{J-1}$
are defined by the deterministic age process, the income process and households'
optimal choices in response to these shocks.} The full numerical computational method is outlined
in detail in the Appendix.

\subsection{Transition Equilibrium and Forecast Rules}

\rev{For a given demographic path and forecast rule indexed by $\lambda$, we
compute a transition path as follows. At each date, households \zac{optimise} using a
perceived continuation path for house prices. Their choices determine voter
composition, voting determines housing-supply restrictions and \zac{the generated price path is the sequence of market-clearing prices conditional on the permitted housing stock.} Let
$\boldsymbol{\Omega}=\{\Omega_{t}\}_{t=0}^{T}$ denote the demographic path and
let $\widehat{\mathcal{P}}^{(n)}=\{\widehat{\mathbf{p}}_{t}^{(n)}\}_{t=0}^{T}$
denote the collection of date-specific continuation price paths used after
forecast revision $n$. Given $\boldsymbol{\Omega}$ and
$\widehat{\mathcal{P}}^{(n)}$, we solve household decisions backward and
simulate the economy forward, applying the political housing-supply rule at
each date. The resulting voter composition and supply restrictions imply a
generated house-price path
$\mathcal{T}\left(\widehat{\mathcal{P}}^{(n)};\boldsymbol{\Omega}\right)$. Let
$\mathcal{T}_{t}\left(\widehat{\mathcal{P}}^{(n)};\boldsymbol{\Omega}\right)$
denote the continuation segment of that generated price path beginning at date
$t$. The bounded-learning specifications start from the current-price rule,
$\widehat{\mathcal{P}}^{(0)}=\widehat{\mathcal{P}}^{C}$, and move each
perceived continuation path part of the way toward the path generated by future
demographics, voting outcomes and \zac{market clearing conditional on permitted supply}:}
\begin{equation}
\revmath{\log \widehat{\mathbf{p}}_{t}^{(n+1)}=\left(1-\lambda\right)\log \widehat{\mathbf{p}}_{t}^{(n)}+\lambda\log \mathcal{T}_{t}\left(\widehat{\mathcal{P}}^{(n)};\boldsymbol{\Omega}\right),\quad t=0,\ldots,T,}
\label{eq:forecastupdate}
\end{equation}
\rev{where $n$ indexes belief revisions rather than calendar time. The
parameter $\lambda\in[0,1]$ indexes the weight placed in each revision on the
path generated by future demographics, voting outcomes and \zac{market clearing conditional on permitted supply}.
When $\lambda=0$, \eqnlink{eq:forecastupdate} leaves perceived future
prices at the current-price rule. Positive values of $\lambda$ move the
perceived path toward the path generated by demographics, voting outcomes,
supply restrictions and \zac{market clearing conditional on permitted supply}. These paths are
finite-revision expectation-sensitivity exercises: they vary how much of the
demographic-political path enters household forecasts rather than selecting a
single expectations concept. We do not interpret $\lambda$ as a numerical
relaxation parameter that is iterated to a common fixed point. If the update
were iterated fully to convergence, positive values of $\lambda$ would target
the same fixed-point object. The reported exercises instead use finite forecast
revisions and compare the resulting subjective forecast paths.}

\rev{The parameter $\lambda$ has a direct numerical interpretation. A forecast
revision means applying \eqnlink{eq:forecastupdate} once to the whole future
price path. If the generated path were held fixed, the cumulative weight placed
on that path after $k$ revisions would be $1-\left(1-\lambda\right)^{k}$. This
gives a useful scale for the plotted values: after ten revisions the
corresponding weight is about 40 per cent for $\lambda=5$ per cent, 65 per cent
for $\lambda=10$ per cent and 89 per cent for $\lambda=20$ per cent.}

\rev{A full rational-expectations transition would instead impose the fixed
point}
\begin{equation}
\revmath{\widehat{\mathbf{p}}_{t}=\mathcal{T}_{t}\left(\widehat{\mathcal{P}};\boldsymbol{\Omega}\right)\quad\text{for all }t,}
\label{eq:full_re_fixed_point}
\end{equation}
\rev{so that households' perceived price path coincides with the path generated
by their own housing choices, voting outcomes, supply restrictions and \zac{market clearing conditional on permitted supply}. This is a strong informational requirement in the present
environment: households would need to forecast not only future prices, but
also how demographic change affects voter composition, political restrictions
on supply, housing choices and \zac{the market-clearing prices implied by permitted supply}.
Our reported exercises
therefore focus on how outcomes vary with the information households use when
forecasting future prices.\footnote{\rev{This interpretation is close to \paperlink{moll2026}{Moll (2026)}, who argues that rational expectations over equilibrium prices in
heterogeneous-agent models place a strong forecasting burden on agents and
motivates restricted forecasting rules in which agents forecast prices more
directly. We implement this idea through bounded-learning forecast rules for
future house prices, with $\lambda$ controlling how much information about the
future demographic and political path enters those forecasts.}}}

\section{Calibration}

We calibrate the model to capture salient features of the US housing
market in the early 2000s prior to the housing boom and bust. The
model period is one year. \rev{Households are economically active from ages
25 through 80, giving $J=56$ annual periods. They work for 41 periods, through
age 65, and receive the terminal bequest payoff at the end of life.} The risk aversion parameter
is set to 2, as is standard in the macroeconomics literature. The
income process consists of the parameters for the deterministic age-profile,
the persistent AR(1) component, and the transitory i.i.d. component
of income. We estimate the parameters of the deterministic and stochastic
income processes using data from the Panel Study of Income Dynamics. 

The retirement income replacement rate is set at 50\% of the final
period's non-transitory income, based on \paperlink{diaz_luengo_prado2008}{D\'iaz and Luengo-Prado (2008)}. The risk-free
interest rate is set at 2\%, while the mortgage interest rate is 5\%.
The maximum loan-to-value (LTV) is set at 0.9 in line with \paperlink{greenwald2018}{Greenwald (2018)}. The spread between the saving rate and the mortgage rate is
set at 3\% matching the 30 year average difference between the two.

\begin{table}[H]
\centering
\scriptsize
\caption{Parameter Values}
\label{tab:parameters}
\vspace{0.35em}
\makebox[\textwidth][c]{%
\begin{tabular}{cc>{\centering\arraybackslash}p{3cm}>{\centering\arraybackslash}p{3cm}}
\toprule 
Parameter & Value & Interpretation & Source\tabularnewline
\midrule
$J$ & 56 & Periods of life & \paperlink{kaplan_mitman_violante2020}{KMV (2020)}\tabularnewline
$J^{retired}$ & 41 & Periods of working life & \paperlink{kaplan_mitman_violante2020}{KMV (2020)}\tabularnewline
$r^{a}$ & 0.02 & Risk free rate & 30 year average\tabularnewline
$r^{m}$ & 0.05 & Mortgage interest rate & 30 year average\tabularnewline
$\chi$ & 0.03 & Mortgage spread & 30 year average\tabularnewline
$\beta$ & $0.978$ & Discount factor & Internally Calibrated\tabularnewline
$\gamma$ & 2 & Risk aversion & \paperlink{kaplan_mitman_violante2020}{KMV (2020)}\tabularnewline
$\zeta$ & 0.75 & Housing weight in utility & Internally Calibrated\tabularnewline
$F^{sell}$ & 0.07 & Cost of selling house & \paperlink{kaplan_mitman_violante2020}{KMV (2020)}\tabularnewline
$m$ & 0.9 & Maximum LTV & \paperlink{greenwald2018}{Greenwald (2018)}\tabularnewline
$\psi$ & 10 & Bequest scaling & Internally Calibrated\tabularnewline
$\overline{w}$ & $150$ & Bequest parameter & Internally Calibrated\tabularnewline
$\sigma$ & $-2.4$ & Pro-supply Non-financial preference & Internally Calibrated, shifts 15\% of the voters in steady
state\tabularnewline
$\phi^{rent}$ & $0.05$ & Rental per unit cost & Internally Calibrated\tabularnewline
$\theta^{rent}$ & 0.9 & Rental discount & Mian et al (2017)\tabularnewline
\bottomrule
\end{tabular}}
\end{table}

\rev{We set the internally calibrated parameters
$\left\{ \beta,\sigma,\zeta,\overline{w},\psi,\phi^{rent}\right\} $
using simulated method moments to match the following moments. The rental
discount $\theta^{rent}$ is taken from the external calibration reported in
Table~\ref{tab:parameters}.}

\begin{table}[H]
\caption{Calibration}
\label{tab:calibration}
\vspace{0.35em}

\centering{}%
\begin{tabular}{c>{\centering}p{2cm}>{\centering}p{2cm}}
\toprule 
 & Model & Data\tabularnewline
\midrule
Share of owners with a mortgage & 0.55 & 0.57\tabularnewline
Share of renters & 0.38 & 0.38\tabularnewline
House price to Income Ratio & 6.0 & 6.0\tabularnewline
Home-ownership rate 25-35 & 0.38 & 0.40\tabularnewline
Home-ownership rate 45-55 & 0.72 & 0.70\tabularnewline
Home-ownership rate 65-75 & 0.71 & 0.75\tabularnewline
\bottomrule
\end{tabular}
\end{table}

The price point at which the median voter is indifferent to additional
housing supply will be affected by $\sigma$, which captures households'
non-financial preferences over the supply of housing. We calibrate
$\sigma$ such that the housing market is in equilibrium in 2000.
Our calibration requires a non-financial preference in favour of approving
more homes as homeowners almost always make up a large majority of
the voting population.

\section{Quantitative Analysis}

\subsection{Steady State}

First we \zac{analyse} how changes in the age distribution affect the housing market
in a set of stationary exercises. \rev{For each year, we take the observed US
age distribution as given and solve for the political-economic equilibrium
implied by the model.} We calculate the political-economic
equilibrium using demographic data from the US over the period 1950-2020.
During this time there was substantial variation in the demographic
structure, driven by the generation born in the post-war spike and
how it reverberated over time, as shown in \figlinklower{fig:3}
at the beginning of the paper. For example, the mean age starts at
30 in 1950 and eventually falls as those born during the baby boom
increase the share of young people in the labour and housing market.
The mean age in the labour market troughs at 28 in 1970 as the baby
boom subsides. The mean age then continues to steadily rise to 39
in 2020 as those born during the baby boom age.

To estimate the effect of the changing demographic composition for
each year in this period, we solve for the equilibrium price that
would result in the median voter being indifferent over an increase
or decrease in housing supply. All other parameters are kept constant
with the only change coming via variation in demographic structure.
\rev{Because each calculation is stationary, the forecast-rule distinction
does not materially enter this exercise; it becomes substantive in the
transition analysis, where voter composition and political support for
housing-supply restrictions change over time.}
Operationally, for each target age distribution we set household formation
so that the demographic structure is constant in the absence of further
shocks, and then solve for the associated stationary political-economic
equilibrium house price.

Since home-ownership rates increase with age, we find that when the
median population age increases the aggregate home-ownership rate rises.
This increases political support for a reduction in the housing supply
which increases house prices. House prices rise until there is enough
political support for additional housing supply (via an increase in
households who rent) such that the median voter is again indifferent
to a small change in house prices.

\figlink{fig:5} shows the model's political-economic equilibrium
cost of housing over the post-war period alongside data on real rents.\footnote{In our model's environment real rents move in line with the cost of
housing, while abstracting from the impact of interest rates which
will impact on the empirical price of housing.} This equilibrium cost represents the effects of demographic shocks
on the voter coalition in the housing market, and thus ignores
other factors such as changes in the business cycle, technology and
financial crises. Accordingly one should not expect it to closely
track actual costs that are affected by these non-demographic factors.
Despite this, the equilibrium results fit the post-war trend
in real rents quite closely.\footnote{We use rents instead of house prices as a measure of housing costs
to avoid the measure being distorted by the secular decline in interest
rates over this period. Real rents are calculated using the average
rent of a primary residence across US cities divided by headline CPI.} In the model the post-war spike in the birth rate leads to a significant
fall in the equilibrium house price starting in the 1970s as these
households, who are mostly young renters, enter the housing market.
As these households age and become homeowners they shift towards opposing
the construction of new housing to protect the value of their stock
of housing. Real rents show a similar trend peaking slightly earlier
in the 1950s, before trending downwards for the next three decades
reaching a low point in the 1980s then increasing steadily up until
2020.
\begin{figure}[H]
\begin{centering}
\includegraphics[width=0.66\textwidth]{Figures/figure5_python_steady_state_house_price.pdf}
\par\end{centering}
\noindent\raggedright{}\caption{\rev{Stationary Accounting House Price and Real Rents}\emph{ }}
\emph{\label{fig:5}}
\end{figure}

Using demographic projections out to 2100 we can estimate the impact
the ongoing ageing of the US economy will have on the political economy
of the US housing market as shown in \figlinklower{fig:6}. \rev{The projection
series report representative-locality demographic paths for the US population under low, medium and
high immigration assumptions. In the low and medium cases, ageing raises the
share of voters who favour tighter supply at the initial price, so the stationary
accounting price rises. Under high immigration, younger inflows partly offset
ageing at first, but those cohorts later age into the homeowner coalition.}

\begin{figure}[H]
\begin{centering}
\includegraphics[width=0.66\textwidth]{Figures/figure6_python_steady_state_projection.pdf}
\par\end{centering}
\noindent\raggedright{}\caption{\rev{Stationary Accounting House Price Projections}}
\emph{\label{fig:6}}
\end{figure}


\subsection{Impulse Response Functions}

We next calculate the dynamics of a demographic shock
to the housing market in the form of a temporary baby boom. We start
from the initial equilibrium and then increase the birth rate of new
households by 10 per cent for 10 years.
After 25 years, when those who were born during the boom enter the
labour and housing market, the effect of the spike in the birth rate
starts to affect the demographic structure of households in the model.
We then calculate forecast-rule transition paths across the
subsequent 80 years to trace the effect of the baby boom through time.

\rev{\figlink{fig:7} reports the perceived house-price paths in the
temporary baby-boom experiment, measured as log deviations from the initial
price. We plot $\lambda=0$, 5, 10 and 20 per cent. When
$\lambda=0$, households do not use demographic-political forecasts. Instead,
future housing conditions in continuation values are \zac{summarised} by the current
house price at date $t$. As $\lambda$ increases, perceived future prices place
more weight on the path implied by future demographics, voting outcomes,
supply restrictions and \zac{market clearing conditional on permitted supply}.}

\rev{The shape of the figure follows the life-cycle timing of the shock.
Perceived prices are nearly flat at first, before rising as the boom generation reaches
housing-market age and changes the relative size of the cohorts whose housing
choices and votes determine supply restrictions. The price path then falls back
because the demographic shock is temporary and the pressure passes through the
age distribution. Higher values of $\lambda$ dampen the hump because
households' forecasts incorporate more of this future demographic and political
adjustment, rather than extrapolating only from current market conditions.}
\begin{figure}[H]
\centering
\includegraphics[width=0.66\textwidth]{Figures/figure7_unified_temporary_babyboom_forecast_rules.pdf}
\caption{Forecast Price Paths for a Temporary Birthrate Increase}
\label{fig:7}
\end{figure}

\rev{\figlink{fig:8} reports net worth and consumption for the generation
born during the temporary baby boom. Each series is measured relative to
households of the same age in the initial stationary economy, so a value above
one means that the boom cohort has higher net worth or consumption than
same-age households before the shock. As in \figlink{fig:7}, higher values
of $\lambda$ dampen the house-price response to the demographic shock, which
compresses the net-worth swing for the boom cohort. Consumption follows a
smoother life-cycle profile, but it first dips below and then rises above the
age-specific initial value as housing and saving decisions adjust over time.}

\begin{figure}[H]
\centering
\includegraphics[width=0.70\textwidth]{Figures/figure8_unified_boom_cohort_forecast_rules.pdf}
\caption{Boom-Cohort Effects of a Temporary Birthrate Increase}
\label{fig:8}
\end{figure}

\section{Robustness and Extensions}

\rev{We report two sensitivity checks in the main text. The first varies the
non-financial preference term that shifts the voting threshold in the temporary
baby-boom transition. The second changes the weights assigned to voters in the
stationary historical exercise. Both checks point to the same interpretation:
changes in voter age and tenure composition drive the results. The magnitude
varies with the political calibration and voting weights, but the demographic
timing of the mechanism is unchanged.}

\subsection{Non-Financial Preference Calibration}

\rev{\figlink{fig:9} varies the non-financial preference term in the voting
rule. This term shifts the level of support for restricting housing supply and
therefore changes the stationary house-price-to-income ratio. The transition
paths are similar. The initial fall, later rise as the boom cohort enters the
housing market, and return toward the stationary allocation change little
because the experiment is driven mainly by voter-composition changes over time.
Changing the non-financial preference term shifts the average level of political
support, but it has little effect on the marginal political force created by the
larger cohort.}

\begin{figure}[H]
\begin{centering}
\includegraphics[width=0.66\textwidth]{Figures/figure9_python_nonfinancial_pref_robust.pdf}
\par\end{centering}
\noindent\raggedright{}\caption{Robustness to Non-Financial Preference Calibration}
\label{fig:9}
\end{figure}


\subsection{Alternative Voting Weights}

\rev{\figlink{fig:10} changes how household votes are aggregated in the
stationary historical exercise. We compare the one-person-one-vote specification
with voting weights based on age-specific turnout and log household net worth.
Turnout weighting has little effect on the implied house-price path because it
changes the political weight of the age groups driving
the demographic mechanism only slightly. Log-net-worth weighting has a larger quantitative effect,
since political influence shifts toward older and wealthier homeowners as the
economy ages. Even in this case, however, the historical increase remains
positive, although smaller.}

\begin{figure}[H]
\begin{centering}
\includegraphics[width=0.66\textwidth]{Figures/figure10_python_alternative_voting_weights.pdf}
\par\end{centering}
\noindent\raggedright{}\caption{Alternative Voting Weights in the Stationary Historical Exercise}
\label{fig:10}
\end{figure}

\rev{Table~\ref{tab:historical_accounting_decomposition} \zac{summarises}
the same historical accounting exercise across the voting weights used
in \figlinklower{fig:10}, together with an additional unlogged net-worth
case. Holding the house price fixed at its year-2000 model value, the historical
age distribution raises support for restricting supply by 8.2 percentage points
between 1950 and 2020 under one-person-one-vote. The stationary
accounting-equilibrium price rises by 14.3 per cent over the same period.
Weighting votes by turnout gives a similar change in political support and a
12.9 per cent price increase. Weighting votes by household net worth lowers the
implied price change, but still produces a 7--8 per cent increase. The historical
result is therefore not driven only by the one-person-one-vote specification,
although the voting weights affect the magnitude of the channel.}

\begin{table}[H]
\centering
\caption{Historical Accounting and Voting Weights}
\label{tab:historical_accounting_decomposition}
\begin{tabular}{lccc}
\toprule
 & \multicolumn{2}{c}{Support for restricting supply} & Equilibrium price\tabularnewline
\cmidrule(lr){2-3}\cmidrule(lr){4-4}
Specification & 1950 & 2020 & 1950--2020 change\tabularnewline
\midrule
One person, one vote & 40.2 & 48.4 & 14.3\%\tabularnewline
Turnout weighted & 43.0 & 51.3 & 12.9\%\tabularnewline
Log net-worth weighted & 45.1 & 52.3 & 8.4\%\tabularnewline
Net-worth weighted & 64.7 & 69.3 & 7.1\%\tabularnewline
\bottomrule
\end{tabular}\par
\begin{minipage}{0.88\textwidth}
\vspace{0.4em}
\footnotesize
\emph{Notes}: The support columns report the share of weighted votes
in favour of restricting housing supply, evaluated at the model's
year-2000 house price. The final column reports the percentage change
in the stationary accounting-equilibrium house price between
1950 and 2020 under the same voting rule. All calculations use the
historical age distributions and hold the remaining model parameters
fixed.
\end{minipage}
\end{table}

\section{Conclusion}

Local planning regulations and the political-economic constraints
that bind them have long been defining features of housing markets
across high-income countries. Because housing is a large component of
homeowners' financial portfolios, some households have a direct financial
incentive to protect this value by voting in favour of regulations that restrict
the supply of new housing. In this paper, we show how post-war demographic
changes shifted relative generation sizes and, consequently, political
support for additional housing supply. In the model, this demographic-political
mechanism implies a trough-to-peak increase in the house-price-to-income ratio
of roughly 15 per cent. Furthermore, future demographic projections imply
continued upward pressure on house prices over the coming decades.

Unpacking the distributional impacts of this generational shift, we
find that a one-off increase in the birth rate changes the median voter over
the life cycle. When the cohort is young, political support shifts toward
additional construction; when it ages into home ownership, political support
shifts toward restrictions on new supply. The transition exercises show that
this mechanism is sensitive to how households form expectations about future
demographic and political conditions. Under the current-price rule, households
do not separately forecast the future demographic-political path, which
amplifies the temporary baby-boom price response. When forecasts place more
weight on that future path, the price hump and the boom cohort's net-worth
swing are dampened. The sensitivity checks indicate that the timing of this
mechanism is not driven by a particular calibration of non-financial voting
preferences or by the main voting weights, although these assumptions can
change the level of the response.

The results also speak to current debates over housing-supply reform. If
demographic change raises the political weight of households with an interest
in preserving housing scarcity, then reforms that expand supply may face
stronger opposition precisely when affordability pressures are rising.
Policies that reduce the private capital losses from new construction, or
that broaden participation in local housing decisions, may therefore be
important complements to planning reform.

\section*{References}

\bibtarget{ahlfeldt2011}{Ahlfeldt, G. M. (2011).} Blessing or curse? Appreciation, amenities
and resistance to urban renewal. Regional Science and Urban Economics,
41(1):32\textendash 45.

\bibtarget{ahlfeldt_maennig2015}{Ahlfeldt, G. M. and Maennig, W. (2015).} Homevoters vs. leasevoters:
A spatial analysis of airport effects. Journal of Urban Economics,
87:85\textendash 99

\bibtarget{albouy_etal2019}{Albouy, D., Behrens, K., Robert-Nicoud, F., \& Seegert, N. (2019).}
The optimal distribution of population across cities. Journal of Urban
Economics, 110, 102-113.

\bibtarget{anagol_etal2021}{Anagol, S., Ferreira, F. V., \& Rexer, J. M. (2021).} Estimating the
economic value of zoning reform (No. w29440). National Bureau of Economic
Research.

\bibtarget{attanasio_etal2012}{Attanasio, O. P., Bottazzi, R., Low, H. W., Nesheim, L., \& Wakefield,
M. (2012).} Modelling the demand for housing over the life cycle. Review
of Economic Dynamics, 15(1), 1-18.

\bibtarget{alesina_rodrik1994}{Alesina, A., \& Rodrik, D. (1994).} Distributive politics and economic
growth. The quarterly journal of economics, 109(2), 465-490.

\bibtarget{bosetti_sims2016}{Bosetti, N., and Sims, S. (2016).} \textquotedblleft{}Stopped: Why people
oppose development in their back yard\textquotedblright{}. Centre for London.

\bibtarget{bunten2017}{Bunten, D. (2017).} Is the rent too high? Aggregate implications of
local land-use regulation. Finance and Economics Discussion Series 2017-064,
Board of Governors of the Federal Reserve System.

\bibtarget{chatterjee_eyigungor2017}{Chatterjee, S., \& Eyigungor, B. (2017).} A tractable city model for
aggregative analysis. International economic review, 58(1), 127-155.

\bibtarget{cheung_meltzer2013}{Cheung, R. and Meltzer, R. (2013).} Homeowners associations and the
demand for local land use regulation. Journal of Regional Science,
53(3):511\textendash 534.

\bibtarget{dehring_depken_ward2008}{Dehring, C. A., Depken II, C. A., and Ward, M. R. (2008).} A direct
test of the homevoter hypothesis. Journal of Urban Economics, 64(1):155\textendash 170.

\bibtarget{diaz_luengo_prado2008}{D\'iaz, A., \& Luengo-Prado, M. J. (2008).} On the user cost and home
ownership. Review of Economic Dynamics, 11(3), 584-613.

\bibtarget{duranton_kerr2015}{Duranton, G., \& Kerr, W. R. (2015).} The logic of agglomeration (No.
w21452). National Bureau of Economic Research.

\bibtarget{duranton_puga2023}{Duranton, G., \& Puga, D. (2023).} Urban growth and its aggregate implications.
Econometrica, 91(6), 2219-2259.

\bibtarget{eggleston_etal2020}{Eggleston, J., Hays, D., Munk, R., \& Sullivan, B. (2020).} The Wealth
of Households, 2017. Washington, DC, USA: US Department of Commerce,
US Census Bureau.

\bibtarget{einstein_etal2019}{Einstein, K. L., Palmer, M., \& Glick, D. M. (2019).} Who participates
in local government? Evidence from meeting minutes. Perspectives on
politics, 17(1), 28-46.

\bibtarget{fernandez_krueger2011}{Fernandez-Villaverde, J., \& Krueger, D. (2011).} Consumption and saving
over the life cycle: How important are consumer durables?. Macroeconomic
dynamics, 15(5), 725-770.

\bibtarget{fischel1999}{Fischel, W. A. (1999).} Does the American way of zoning cause the suburbs
of metropolitan areas to be too spread out?. Governance and opportunity
in metropolitan America, 15, 1-19.

\bibtarget{fischel2001}{Fischel, W. A. (2001).} The homevoter hypothesis: How home values influence
local government taxation, school finance, and land-use policies.
Cambridge, MA: Harvard University Press.

\bibtarget{fischel2004}{Fischel, W. A. (2004).} An economic history of zoning and a cure for its
exclusionary effects. Urban Studies, 41(2), 317-340.

\bibtarget{frieden1979}{Frieden, B. J. (1979).} The environmental protection hustle. Cambridge,
MA: MIT Press.

\bibtarget{fu2019}{Fu, C. H. (2019).} Living arrangement and caregiving expectation: The
effect of residential proximity on inter vivos transfer. Journal of
Population Economics, 32(1), 247-275.

\bibtarget{glaeser_gyourko_saks2006}{Glaeser, E. L., Gyourko, J., \& Saks, R. E. (2006).} Urban growth and
housing supply. Journal of economic geography, 6(1), 71-89.

\bibtarget{greenwald2018}{Greenwald, D. (2018).} The mortgage credit channel of macroeconomic
transmission. Research Paper no. 5184-16, Sloan School Management,
Massachusetts Inst. Technology.

\bibtarget{greenwood_etal2005}{Greenwood, J., Seshadri, A., \& Vandenbroucke, G. (2005).} The baby
boom and baby bust. American Economic Review, 95(1), 183-207.

\bibtarget{gyourko_mayer_sinai2013}{Gyourko, J., Mayer, C., \& Sinai, T. (2013).} Superstar cities. American
Economic Journal: Economic Policy, 5(4), 167-199.

\bibtarget{gyourko_molloy2014}{Gyourko, J., \& Molloy, R. (2014).} Regulation and housing supply
(No. w20536). National Bureau of Economic Research.

\bibtarget{hankinson2018}{Hankinson, M. (2018).} When do renters behave like homeowners? High
rent, price anxiety, and NIMBYism. American Political Science Review,
112(3), 473-493.

\bibtarget{herkenhoff_ohanian_prescott2018}{Herkenhoff, K. F., Ohanian, L. E., \& Prescott, E. C. (2018).} Tarnishing
the golden and empire states: Land-use restrictions and the US economic
slowdown. Journal of Monetary Economics, 93, 89-109.

\bibtarget{hilber_robert_nicoud2013}{Hilber, C. A. and Robert-Nicoud, F. (2013).} On the origins of land
use regulations: Theory and evidence from US metro areas. Journal
of Urban Economics, 75:29\textendash 43.

\bibtarget{iacoviello2005}{Iacoviello, M. (2005).} House prices, borrowing constraints, and monetary
policy in the business cycle. American economic review, 95(3), 739-764.

\bibtarget{iacoviello_pavan2013}{Iacoviello, M., \& Pavan, M. (2013).} Housing and debt over the life
cycle and over the business cycle. Journal of Monetary Economics,
60(2), 221-238.

\bibtarget{jiang2018}{Jiang, B. (2018).} Homeownership and voter turnout in US local elections.
Journal of Housing Economics, 41, 168-183.

\bibtarget{justiniano_etal2015}{Justiniano, A., Primiceri, G. E., \& Tambalotti, A. (2015).} Household
leveraging and deleveraging. Review of Economic Dynamics, 18(1), 3-20.

\bibtarget{kaplan_mitman_violante2020}{Kaplan, G., Mitman, K., \& Violante, G. L. (2020).} The housing boom
and bust: Model meets evidence. Journal of Political Economy, 128(9),
3285-3345.

\bibtarget{kahn_vaughn_zasloff2010}{Kahn, M. E., Vaughn, R., \& Zasloff, J. (2010).} The housing market
effects of discrete land use regulations: Evidence from the California
coastal boundary zone. Journal of Housing Economics, 19(4), 269-279.

\bibtarget{kogan_etal2018}{Kogan, V., Lavertu, S., \& Peskowitz, Z. (2018).} Election timing,
electorate composition, and policy outcomes: Evidence from school districts.
American Journal of Political Science, 62(3), 637-651.

\bibtarget{kindermann_krueger2022}{Kindermann, F., \& Krueger, D. (2022).} High marginal tax rates on
the top 1 percent? Lessons from a life-cycle model with idiosyncratic
income risk. American Economic Journal: Macroeconomics, 14(2), 319-366.

\bibtarget{krusell_smith1998}{Krusell, P., \& Smith, A. A., Jr. (1998).} Income and wealth heterogeneity
in the macroeconomy. Journal of Political Economy, 106(5), 867\textendash 896.

\bibtarget{kulka_etal2022}{Kulka, A., Sood, A., \& Chiumenti, N. (2022).} How to increase housing
affordability: Understanding local deterrents to building multifamily
housing (No. 22-10). Working Papers.

\bibtarget{li2022}{Li, X. (2022).} Do new housing units in your backyard raise your rents?
Journal of Economic Geography, 22(6), 1309-1352.

\bibtarget{mankiw_weil1989}{Mankiw, N. G., \& Weil, D. N. (1989).} The baby boom, the baby bust,
and the housing market. Regional science and urban economics, 19(2),
235-258.

\bibtarget{mast2024}{Mast, E. (2024).} Warding off development: Local control, housing supply,
and nimbys. Review of Economics and Statistics, 106(3), 671-680.

\bibtarget{manville_etal2020}{Manville, M., Monkkonen, P., \& Lens, M. (2020).} It's
time to end single-family zoning. Journal of the American Planning
Association, 86(1), 106-112.

\bibtarget{moll2026}{Moll, B. (2026).} The trouble with rational expectations in heterogeneous
agent models: A challenge for macroeconomics. The Economic Journal,
136(676), 1173-1216.

\bibtarget{mumtaz_sustek2023}{Mumtaz, H. and \v{S}ustek, R. (2023).} Global house prices since 1950. Working
Paper.

\bibtarget{oliver_ha2007}{Oliver, J. E., \& Ha, S. E. (2007).} Vote choice in suburban elections.
American Political Science Review, 101(3), 393-408.

\bibtarget{ortalo_magne_prat2014}{Ortalo-Magn\'e, F., \& Prat, A. (2014).} On the political economy of
urban growth: Homeownership versus affordability. American Economic
Journal: Microeconomics, 6(1), 154-181.

\bibtarget{parkhomenko2023}{Parkhomenko, A. (2023).} Local causes and aggregate implications of
land use regulation. Journal of Urban Economics, 138, 103605.

\bibtarget{persson_tabellini1994}{Persson, T., \& Tabellini, G. (1994).} Does centralization increase
the size of government?. European Economic Review, 38(3-4), 765-773.

\bibtarget{pennington2021}{Pennington, K. (2021).} Does building new housing cause displacement?
The supply and demand effects of construction in San Francisco.

\bibtarget{sahn2024}{Sahn, A. (2024).} Racial diversity and exclusionary zoning: Evidence
from the great migration. Journal of Politics. 

\bibtarget{scally_tighe2015}{Scally, C. P., \& Tighe, J. R. (2015).} Democracy in action?: NIMBY
as impediment to equitable affordable housing siting. Housing Studies,
30(5), 749-769.

\bibtarget{song2021}{Song, J. (2021).} The effects of residential zoning in US housing markets.
Available at SSRN 3996483.

\bibtarget{titman_zhu2024}{Titman, S., \& Zhu, G. (2024).} City characteristics, land prices and
volatility. Journal of Urban Economics, 140, 103645.

\bibtarget{trounstine2020}{Trounstine, J. (2020).} The geography of inequality: How land use regulation
produces segregation. American Political Science Review, 114(2), 443-455.

\bibtarget{turner_etal2014}{Turner, M. A., Haughwout, A., \& van der Klaauw, W. (2014).} Land use
regulation and welfare. Econometrica, 82(4), 1341-1403.

\bibtarget{whittemore_bendor2018}{Whittemore, A. H., \& BenDor, T. K. (2018).} Talking about density:
An empirical investigation of framing. Land use policy, 72, 181-191.

\bibtarget{yang2009}{Yang, F. (2009).} Consumption over the life cycle: How different is
housing?. Review of Economic Dynamics, 12(3), 423-443.

\newpage{}

\appendix

\section{Numerical Algorithm}

\rev{The numerical solution has two parts. The first solves the stationary
political-economic accounting equilibrium for a fixed demographic structure.
The second solves transition paths under a specified rule for expectations.}

\subsection{Stationary Equilibrium}

For a fixed demographic structure, the numerical solution follows these steps.
\begin{enumerate}
\item The household's value functions are solved recursively for a fixed
value of $p^{h}$, starting in the final period of life and working
backwards. This provides solutions for
the household value functions $\left\{ V^{rent}_{j},V^{adjust}_{j},V^{non-adjust}_{j}\right\} $
and the resulting decision rules for consumption, housing and liquid
assets as a function of the households' idiosyncratic state variables.
\item The rental price of housing is pinned down by the cost of housing
$p^{h}$.
\item We then calculate how these value functions change if housing supply
is marginally restricted. Operationally, this is evaluated through the
small persistent price change $\epsilon^{p}$ implied by lower supply. The
difference in value functions is used to calculate preferences over housing
supply across households, $\Theta$.
\item Next we solve for the density of households across the life cycle
using these decision rules to solve for the distribution of households
across the idiosyncratic state variables.
\item \rev{We then use these solutions to calculate the share of households
voting to restrict housing supply, $S(p^{h})$. We solve for the root of
$S(p^{h})-0.5=0$ using a bracketed price search. This avoids relying on a
directional update rule and directly implements the median-voter condition.}
\item \rev{Given fixed age weights and household formation rates, we iterate these
decision rules until the cross-sectional distribution is stationary under the
imposed demographic structure.}
\end{enumerate}

\subsection{Transition Paths with Forecast Rules}

\rev{The main text defines the forecast path and the update rule in equations
\reflink{eq:forecastpath} and \reflink{eq:forecastupdate}. The numerical \zac{transition path}
calculation implements these objects as follows.}
\begin{enumerate}
\item \rev{Fix the demographic path
$\boldsymbol{\Omega}=\{\Omega_{t}\}_{t=0}^{T}$ and an initial collection of
perceived continuation price paths. For the current-price case, future prices
are set equal to the current house price at date $t$.}
\item \rev{Given the perceived continuation prices, solve the household problem
backwards along the transition using the same tenure, saving and housing choice
problem as in the stationary case.}
\item \rev{Simulate the economy forward using the resulting decision rules. In
each period, aggregate the voting rule and determine the housing-supply restriction.
\zac{Conditional on the permitted housing stock, compute the market-clearing house price.}
This gives the generated price path $\mathcal{T}\left(\widehat{\mathcal{P}};\boldsymbol{\Omega}\right)$.}
\item \rev{For $\lambda>0$, update the perceived continuation paths toward the generated
path using \eqnlink{eq:forecastupdate}. The reported transition exercises use
a finite number of forecast revisions and report the associated revision counts
and forecast discrepancies in the numerical output.}
\end{enumerate}

\end{document}
